\documentstyle[%


righttag%


,11pt%


,amssymb%


,russian%               remove in Eng.


%,izvru%                 Set izven% in Eng.


,izvru-ks%             for brief communications


]{amsart}





% 





\newcommand{\tts}[1]{}





%\hoffset=-15mm%                  For Eng.


 


\setcounter{page}{73}


\god{1999}


\nomer{9~(448)} %                        Remove in Eng.


%\nomer{Vol.\,43, No.\,9}%               Set in Eng.


%\str{\pageref{begin}--\pageref{end}}%  Set in Eng.


\udk{512}





\author{\it М.П.\,Бурлаков}


% NB:   ^^ remove in Eng.


\title{Обобщенные внешние алгебры и их приложения}





\date{}


%\pagestyle{myheadings}%         Set in Eng.


\pagestyle{plain} %              Remove in Eng.


\begin{document}


%\thispagestyle{plain}%          Set in Eng.


\maketitle


\label{begin}





{\bf 1.} Внешние алгебры Грассмана и Клиффорда давно стали  популярным


инструментом в различных областях математики и теоретической физики.


В основе всех приложений внешних алгебр лежит тождество


\begin{equation}


\bold v \land {\bold v} = {\bold 0},


\end{equation}


где ${\bold v}$ --- вектор подстилающего линейного пространства алгебры


Грассмана. В этой статье рассмотрим обобщения внешних алгебр, в которых


выполняется тождество


\begin{equation}


{\bold v}^m \equiv  \underbrace{ {\bold v \cdot\bold v \dotsb\bold v}}_{m}


= {\bold 0}.


\end{equation}


Ясно, что тождество (1) есть частный случай тождества (2) при $m=2$, и


это обобщение может быть использовано для модификаций разнообразных


приложений внешних алгебр.





Следует однако отметить, что существование внешней алгебры с


характеристическим тождеством (2) далеко не очевидно. Есть простое


рассуждение, кочующее из монографии в монографию, и утверждающее, что


единственной внешней алгеброй может быть только алгебра с антикоммутативным


законом умножения. Мы приведем здесь эти рассуждения, чтобы читатель мог


оценить некоторые особенности конструкции обобщенной внешней алгебры,


а также ``субъективные причины'' по которым внешние алгебры с определяющим


тождеством (2) выпали из поля зрения математиков.





Итак, пусть ${\bold E}_n$ --- линейное векторное пространство


с базисом ${\bold e}_i$, удовлетворяющим перестановочным соотношениям


общего вида


\begin{equation}


{\bold e}_i\cdot {\bold e}_j = \alpha {\bold e}_j\cdot {\bold e}_i.


\end{equation}


Тогда


$$


{\bold e}_i\cdot {\bold e}_j = \alpha\, {\bold e}_j\cdot {\bold e}_i =


\alpha^2\, {\bold e}_i\cdot {\bold e}_j,


$$


т.\,е. $\alpha^2=1$ и тогда либо $\alpha = 1,$ либо $\alpha =-1.$ Если


выбрать $\alpha = 1,$ то над ${\bold E}_n$ получим коммутативную алгебру,


изоморфную алгебре полиномов $n$ переменных; если же выбрать $\alpha =-1,$


то получается внешняя алгебра Грассмана.


\medskip





{\bf 2.} Это простое рассуждение, на первый взгляд, исключает саму возможность


существования внешних алгебр с определяющими тождествами вида (3). Однако


мы можем обойти это препятствие, если перестановочные соотношения (3)


подчиним некоторой упорядоченности в базисе пространства ${\bold E}_n$,


например, положив


\begin{equation}


\begin{split}


{\bold e}_i\cdot {\bold e}_j = \alpha\, &{\bold e}_j\cdot {\bold e}_i,


\quad i>j;\\


{\bold e}_i^m &= 0,


\end{split}


\end{equation}


где $\alpha$ --- примитивный корень $m$-й степени из единицы, а само линейное


пространство ${\bold E}_n$ рассматривается над полем комплексных чисел.





Из формул (4) следуют некоторые простейшие свойства обобщенных внешних


алгебр, которые будем обозначать


${\bold B}_n^m$ или ${\bold B}^m({\bold E}_n)$.


Прежде всего отметим, что если $p<q$, то


\begin{equation}


{\bold e}_p\cdot {\bold e}_q = \alpha^{m-1}\,{\bold e}_q\cdot {\bold e}_p


\equiv \ov{\alpha}\, {\bold e}_q\cdot {\bold e}_p.


\end{equation}


Далее очевидно, что общий член алгебры ${\bold B}_n^m $ представляет собой


неоднородную сумму мономов вида


$\lambda {\bold e}_1^{k_1}\cdot{\bold e}_2^{k_2}


\dotsb{\bold e}_n^{k_n}$, где $0\le k_i <m$


и ${\bold e}^0 = 1$, т.\,е.


\begin{equation}


{\bold z} = \sum_{0\le k_i <m} z_{k_1,k_2,\dotsc, k_n}


{\bold e}_1^{k_1}\cdot {\bold e}_2^{k_2}\dotsb{\bold e}_n^{k_n},


\end{equation}


если $\bold z\in\bold B^m_n$.


Из общего вида элементов обобщенной внешней алгебры следует, что


$\dim_{\bold c} {\bold B}_n^m = m^n.$





Пусть $\wt{\bold B}_n^m$ --- подалгебра алгебры ${\bold B}_n^m$, в которой


отсутствуют ``скалярные'' слагаемые элементов; очевидно $\dim_{\bold c}


\wt{\bold B}_n^m = m^n - 1$ и алгебра $\wt {\bold B}_n^m$ градуирована


$$


\wt {\bold B}_n^m = \,^1{\bold B}_n^m \oplus\, ^2{\bold B}_n^m


\oplus \dotsb \oplus \, ^N{\bold B}_n^m,


$$


где $^k{\bold B}_n^m $ представляет собой линейное пространство сумм


мономов вида $\lambda {\bold e}_1^{k_1}\cdot {\bold e}_2^{k_2}\dotsb


{\bold e}_n^{k_n}$ для которых $k_1+k_2+\dotsb +k_n =k $ и


$ N=(m-1)n$. В силу этого обстоятельства для любых элементов


${\bold z}\in \wt{\bold B}_n^m $ существует степень $M$,


при которой ${\bold z}^M=0.$


В частности, последнее равенство всегда выполняется для $M>N.$


Однако эта общая оценка показателя степени, аннулирующей элементы обобщенной


внешней алгебры, заведомо завышена для частных случаев.





\begin{thm} Пусть ${\bold z}=z_1{\bold e}_1+\dotsb + z_n{\bold e}_n$ ---


вектор


подстилающего пространства ${\bold E}_n$, тогда справедливо тождество


\begin{equation}


{\bold z}^m=(z_1{\bold e}_1+\dotsb +z_n{\bold e}_n)^m=0.


\end{equation}


\end{thm}





Эта теорема в теории обобщенных внешних


алгебр играет центральную роль. При $m=2$,


т.\,е. в случае алгебр Грассмана, тождество (7) выражает антикоммутативность


умножения. Для малых размерностей $n$ и показателей $m$ тождество (7)


легко доказать прямым вычислением; в общем случае доказательство можно


провести индукцией по размерности $n$, предварительно доказав это тождество


в случае $n=2$,


$$


(x{\bold e}_1+ y{\bold e}_2)^m=


P_1^m(\alpha)x^{m-1}y{\bold e}_1^{m-1}{\bold e}_2+


P_2^m(\alpha)x^{m-2}y^2{\bold e}_1^{m-2}{\bold e}_2^2 +


\dotsb + P_{m-1}^m(\alpha)xy^{m-1}{\bold e}_1{\bold e}^{m-1}_2=0.


$$


Доказательство этого последнего тождества опирается на перестановочные


соотношения в обобщенной внешней алгебре.


\medskip





{\bf 3.} Приложения обобщенных внешних алгебр зависят от реализации


их конструкций на различных объектах и поэтому столь же


разнообразны как и приложения внешней алгебры Грассмана.





Отметим, во-первых, приложения обобщенных внешних алгебр к теории


гипердетерминантов. В градуированной алгебре $\wt{\bold B}_n^m$


размерность подпространства $^N{\bold B}_n^m$ равна единице, а следовательно,


произведение $N$ векторов ${\bold z}^{\nu}= z_1^{\nu}{\bold e}_1+\dotsb +


z_n^{\nu}{\bold e}_n$, где $\nu=\overline{1,N}$, дает моном


$$


\prod_{\nu=1}^N = Z{\bold e}_1^{m-1}\cdot {\bold e}_2^{m-1}


\dotsb {\bold e}_n^{m-1}.


$$


Коэффициент этого монома $Z$ представляет относительный инвариант системы


векторов ${\bold z}^{\nu}$ и называется $m$-детерминантом прямоугольной


матрицы, составленной из коэффициентов этих векторов. При $m=2$ коэффициент


$Z$ дает обыкновенный определитель квадратной матрицы.





Существуют многочисленные и разветвленные связи теории $m$-гипердетерминантов


с теорией гипердетерминантов Райта и Кели. Заметим однако, что при помощи


обобщенных внешних алгебр можно получить широкий набор относительных


инвариантов, если вместо векторов подстилающего пространства выбирать другой


набор элементов. Среди этих инвариантов можно выделить один, который


более других походит на обыкновенный детерминант. Для его определения


заметим, что в алгебре $\wt{\bold B}_n^m$ существует подалгебра


$\wt{\bold B}_n$, порожденная элементами ${\bold e}_k^{m-1}=\varepsilon_k$.


Для этих элементов перестановочные соотношения имеют вид


$\varepsilon_i \cdot \varepsilon_j = \alpha \varepsilon_j \cdot


\varepsilon_i$, $j<i$. Но $\varepsilon_i^2 = 0$. Отсюда сразу следует,


что $\dim\wt{\bold B}_n= 2^n-1$, и произведение $n$ элементов этой алгебры


${\bold x}^{\nu}=x_1^{\nu}\varepsilon_1 +\dotsb +x_n^{\nu}\varepsilon_n$


даeт моном максимальной степени


$$


\prod_{k=1}^n {\bold x}^{\nu}= X{\bold e}^{m-1}_1\dotsb


{\bold e}^{m-1}_n\,=\, X\varepsilon_1\dotsb\varepsilon_n.


$$


Коэффициент $X$ называется $m$-детерминантом квадратной матрицы, составленной


из коэффициентов элементов ${\bold x}^k$, и обладает свойствами, аналогичными


свойствам простого детерминанта. Например, если $m<n$ и в матрице


$(x^{\nu}_i)$\; $m$ одинаковых строк, то $m$-детерминант такой матрицы


равен нулю.





Отметим еще, что $m$-детерминанты занимают ``промежуточное'' положение между


обыкновенными определителями, т.\,е. детерминантами и перманентами, и


в силу этого могут использоваться при подсчете числа состояний в


парастатистиках порядка $m$, промежуточных между статистиками Ферми и Бозе.


\medskip





{\bf 4.} Другое приложение обобщенных внешних алгебр находим в теории


дифференциальных форм на комплексных (и почти комплексных) многообразиях.


Пусть ${\bold M_c}^n$ --- комплексное многообразие и рассмотрим на нем


``внешние'' дифференциальные формы вида


$$


\omega = \sum_{0\le k_i <m} \omega_{k_1\dotsc K_n}dz_1^{k_1}


\dotsb dz_n^{k_n},


$$


очевидно, на множестве таких дифференциальных форм реализуется


(бесконечномерная) обобщенная внешняя алгебра, которую обозначим


${\bold B}^n({\bold M_c}^n)$. В этой алгебре существует дифференциальный


оператор, обобщающий оператор внешнего дифференцирования Картана, который


в координатной окрестности определяется по формуле


\begin{equation}


{\bold d}\omega = \sum_{0\le k_i<m} d\omega_{k_1\dotsc k_n}\cdot


dz_1^{k_1}\dotsb dz_n^{k_n}.


\end{equation}


Заметим, что при $m=2$ формула (8) дает оператор внешнего


дифференцирования Картана. Оператор Картана ${\bold d}$ удовлетворяет


тождеству ${\bold d\,\bold d=0}$. Модификацию этого тождества для обобщенной


внешней алгебры дает





\begin{thm} Оператор обобщенного внешнего дифференцирования


удовлетворяет операторному тождеству


$$


{\bold d}^m= \underbrace{{\bold d \dotsb \bold d }}_m= 0.


$$


\end{thm}





Отметим теперь, что действие оператора ${\bold d}^q$ согласовано с


градуировкой алгебры ${\bold B}^m({\bold M_c}^n)$, т.\,е.


$$


{\bold d}^q:\,^p{\bold B}^m({\bold M_c}^n)\to


\,^{p+q}{\bold B}^m({\bold M_c}^n),


$$


и в соответствии с этим дадим следующие определения.





1) Форма $\omega^p\in \,^p{\bold B}^m({\bold M_c}^n)$


называется $q$-циклической


или $q$-циклом, если ${\bold d}^q\,\omega^p=0$.





2) Форма $\omega^p$ называется $q$-точной, если существует форма


$\theta \in \,^{p+q-m}{\bold B}^m({\bold M_c}^n)$ такая, что $\omega^p


={\bold d}^{m-q}\theta$.





Множество $q$-циклов образуeт аддитивную группу (и даже линейное пространство)


${\bold X}^p\subset \,^p{\bold B}^m({\bold M_c}^n);$ множество


$q$-точных $p$-форм


также образуют аддитивную группу (и линейное пространство)


${\bold Y}^p\subset \,^p{\bold B}^m({\bold M_c}^n).$ Легко видеть, что


${\bold Y}^p\subset {\bold X}^p,$ т.\,к. для $q$-точной формы $\omega^p$


имеем


$$


{\bold d}^q\, \omega^p = {\bold d}^q\: {\bold d}^{m-q}\theta = {\bold d}^m


\theta = 0,


$$


и это приводит еще к одному определению.





3) Факторгруппа аддитивной группы $q$-циклов ${\bold X}^p$ по аддитивной


группе $q$-точных $p$-форм ${\bold Y}^p$ будем называть группой


$(p,q)$-гиперкогомологий комплексного


(или почти комплексного) многообразия ${\bold M_c}^n$ и обозначать


${\bold H}_q^p({\bold M_c}^n)$. Это определение обобщает понятие


групп когомологий де Рама, которые соответствуют случаю $m=2$.


\medskip





{\bf 5.} Дуальной к теории когомологий является теория гомологий, в основе


которой лежит понятие граничного оператора. Приведем вариант теории для


групп $k$-мерных цепей над полем комплексных чисел; элементами группы


$k$-мерных


цепей будут формальные линейные комбинации с комплексными коэффициентами


$k$-мерных симплексов, взятых из некоторого симплициального комплекса


${\bold K_c}$, группу $k$-мерных цепей обозначим ${\bold C}_k={\bold C}_k


({\bold K_c})$. Определим граничный оператор  $\delta$, который на симплексе


$ \Delta_k = [A_1,A_2,\dotsc,A_{k+1}] $ действует по правилу


$$


\delta\,\Delta_k = \sum_{i=1}^{k+1}\, \alpha^{i-1}[A_1,\dotsc,A_{i-1},


A_{i+1},\dotsc,A_{k+1}]


$$


и на произвольную $k$-цепь распространяет свое действие по линейности.





\begin{thm} Для любой k-цепи ${\bold c}\in {\bold C}_k({\bold K_c})$


имеет место тождество


$$


\underbrace{\delta\dotsc\delta}_m {\bold c}=\delta^m{\bold c}=0.


$$


\end{thm}





В алгебре Грассмана операторы ${\bold d}$ и ${\boldsymbol \delta}$ называются


соответственно кограничным и граничным, и связь между ними дается общей


интегральной теоремой Стокса--Пуанкаре. Аналогичная теорема имеет место


и в случае обобщенной внешней алгебры.





\begin{thm}


Для формы $\omega$ и цепи ${\bold z}$ имеют место тождества


$$


\int_{(\delta^{m-1} {\bold z})} \omega = \int_{(\delta^{m-2} {\bold z})}


{\bold d}\omega= \dotsb = \int_{(\bold z)} {\bold d}^{m-1} \omega.


$$


\end{thm}





{\bf 6.} Если на линейном пространстве ${\bold E}_n$ задана метрика


$g$,


то алгебра Грассмана может быть перестроена в алгебру Клиффорда


переопределением произведения


$$


{\bold e}_i\cdot {\bold e}_j= {\bold e}_i\wedge {\bold e}_j+ g({\bold e}_i


{\bold e}_j).


$$


Аналогично, если на линейном пространстве задана $m$-форма, то обобщенная


алгебра Грассмана может быть перестроена в обобщенную алгебру Клиффорда


${\bold C}_n^m$


с перестановочными соотношениями (4), (5) и определяющими тождествами


вида ${\bold e}_k^m=1$. Для обобщенных алгебр Клиффорда имеет место теорема,


аналогичная теореме 1.





\begin{thm}


В обобщенной алгебре Клиффорда ${\bold C}_n^m $ для


любого вектора ${\bold z}=z_1{\bold e}_1+\dotsb+z_n{\bold e}_n$ справедливо


тождество


$$


{\bold z}^m= (z_1{\bold e}_1+\dotsb +z_n{\bold e}_n)^m=z_1^m+\dotsb+z^m_n.


$$


\end{thm}





Конструкция обобщенной алгебры Клиффорда порядка $m$ может быть использована


в задаче факторизации линейных дифференциальных уравнений в частных


производных порядка $m$, аналогичной задаче Дирака при факторизации


уравнения Клейна--Гордона--Фока. А именно, имеет место





\begin{thm}


Если функция


$$


{\bold f}= \sum_{0\le k_i<m} f_{k_1\dotsc k_n}\,{\bold e}_1^{k_1}


\dotsb {\bold e}_n^{k_n} :{\bold M}_n\to{\bold C}_n^m


$$


удовлетворяет уравнению


$$


\nabla {\bold f}=0,


$$


где


$$


\nabla = {\bold e}_1\frac{\partial}{\partial z_1}+\dotsb +


{\bold e}_n \frac{\partial}{\partial z_n},


$$


то каждая компонента $f_{k_1\dotsc k_n} $ функции ${\bold f}$ удовлетворяет


уравнению


$$


\frac{\partial^m f_{k_1\dotsc k_n}}{\partial z_1^m}+\dotsb +


\frac{\partial^m f_{k_1\dotsc k_n}}{\partial z_n^m}= 0.


$$


\end{thm}





Неприводимые (относительно действия $\nabla$) элементы алгебры ${\bold C}_n^m$


можно считать аналогами спиноров со значением спина ${\bold s}=\frac{1}{m}$.





\begin{thebibliography}{9}





\bibitem{1}


Бурлаков М.П. {\em Внешние алгебры и их приложения.} --


Грозный, 1991. -- 134 c. 


\bibitem{2}


Бурлаков М.П. {\em Обобщенные алгебры Грассмана на комплексных


многообразиях} // Дифференц. геометрия многообразий фигур. --


Калининград, 1991. -- Вып. 22. -- C.\,14--19.





\end{thebibliography}





\vskip 2cm





\post{\it Московский государственный}{\it Поступила}


\post{\it педагогический университет}{14.04.1999}





\label{end}


\end{document}





